\section{Theoretical Background}
\label{sec:background}

\paragraph{Regression and generalisation.}
We predict the close of session $t+1$ from features known at the close of session
$t$. Training error usually falls as capacity grows, so performance must be judged
on unseen data \citep{hastie2009elements}.

\paragraph{Model families.}
Ridge regression is a regularised linear model. Support vector regression with an
RBF kernel \citep{drucker1997support} predicts from weighted kernel similarities,
so it depends on distances in feature space. A random forest
\citep{breiman2001random} averages bootstrapped trees to reduce variance.
Histogram gradient boosting \citep{friedman2001gradient} fits trees sequentially
to residuals and can also fit noise.

\paragraph{Preprocessing is part of the model.}
A tree splits on $x_j < \tau$. Rescaling $x_j$ by a positive constant preserves
order, so an equivalent threshold exists. A kernel method sums squared differences
\emph{across} features, so rescaling changes which feature dominates. We test this
asymmetry in Section~\ref{sec:ablation}.

\paragraph{Ordered data and leakage.}
Random cross-validation lets a model interpolate across time and is invalid here
\citep{bergmeir2018note,roberts2017crossvalidation}. We use chronological splits
with a boundary purge. Leakage, meaning information unavailable at prediction
time, can make an evaluation optimistic \citep{kaufman2012leakage}, so we test for
it.
